\section{Teoretisk bakgrunn}
\label{sec:background-nn}

Teorien sameinar tre formelle tradisjonar som har utvikla seg uavhengig av kvarandre. Kvar adresserer ein naudsynt aspekt av designform; syntaks, epistemisk logikk og dynamikk; men ingen gjev ein fullstendig modell åleine.

\subsection{Formalgebra og formgrammatikk: syntaksen til forma}

George Stiny og James Gips introduserte formgrammatikkar i 1972 som ein generativ spesifikasjon for visuelt design \parencite{Stiny1972}. Ein formgrammatikk er eit regelsystem som opererer på former gjennom innleiring: ein regel gjenkjenner ei delform i den noverande forma og erstattar ho med ei anna. Kjerneinnsikta er at reglar fyrer på den \textit{noverande geometrien}, ikkje på konstruksjonshistoria.

Stiny formaliserte den underliggjande algebraen \parencite{Stiny1991, Stiny2006}. Former er element i ein boolsk algebra lukka under sum (union), differanse og produkt (snitt). Transformasjonsgruppa til det euklidske rommet gjev rotasjon, translasjon og skalering. Innleiringsrelasjonen $\tau(a) \leq s$ avgjer om ei delform $a$ kan finnast i ei form $s$ under ein transformasjon $\tau$.

Formgrammatikkar gjev ein presis syntaktisk modell: dei spesifiserer \textit{kva operasjonar som er moglege}. Men dei forklarer ikkje \textit{kvifor} éi form oppstår framfor ei anna. Grammatikken genererer eit språk av moglege former; ho forklarer ikkje kvifor somme former vert realiserte og andre ikkje. Dette forklaringsgapet er det fyrste den noverande teorien adresserer.

\subsection{C-K-teori: logikken i design}

Armand Hatchuel og Benoit Weil introduserte C-K-teori (Concept-Knowledge) som ein formell modell av innovativ designtenking \parencite{Hatchuel2003, Hatchuel2009}. Teorien partisjonerer designaren sitt kognitive rom i to utvidbare rom: eit \textit{tankerom} $C$ (proposisjonar med uavgjort sannheitsverdi) og eit \textit{kunnskapsrom} $K$ (proposisjonar med kjend sannheitsverdi). Design skrid fram gjennom fire operatorar som flyttar proposisjonar mellom dei to romma.

C-K-teori har møtt to vedvarande kritikkar relevante for det noverande arbeidet. For det fyrste er tankerommet $C$ i praksis eindimensjonalt: konsept vert spesifiserte gjennom progressiv partisjonering, men rommet manglar geometrisk struktur. For det andre gjev C-K-teori ingen modell av funksjon eller struktur som konstituerande element i designobjektet.

Den noverande teorien adresserer begge. Morforommet gjev $C$ flerdimensjonal geometrisk struktur; konsept er posisjonar i morforommet, ikkje greiner på eit tre. Affordansestrukturen erstattar funksjonsomgrepet med ein relasjonell konstruksjon som er både formelt presis og empirisk testbar.

\subsection{Multiskala-kompetanse: dynamikken til forma}

Michael Levin sitt arbeid med multiskala-kompetanse gjev den tredje søyla: ein substrat-uavhengig modell av målorientert åtferd over biologiske skalaer \parencite{Levin2022}. Levin demonstrerer at agens; definert som målorientert navigasjon via tilbakekopling; opererer på kvart nivå av biologisk organisering.

Chris Fields og Levin formaliserer dette gjennom omgrepet \textbf{kognitiv lyskjegle}: delmengda av omgjevnadene agenten kan representere og manipulere \parencite{Fields2022}. Alt utanfor lyskjegla manglar kausal eksistens for agenten.

Doctor, Wakefield, Pavlic og Levin føreslår \textbf{omsorg}; talet på dimensjonar agenten aktivt representerer; som drivkraft for intelligens \parencite{Doctor2024}. Ein agent som har omsorg for fleire aksar navigerer eit rikare landskap og finn løysingar som toler fleire trykk samstundes. Intelligens er ikkje prosesseringshastigheit; det er dimensjonaliteten til engasjert merksemd.

Dette rammeverket har ikkje tidlegare vore nytta i designteori. Likevel er kartlegginga naturleg. Designaren er ein agent som navigerer eit morforomslandskap via grammatikkoperasjonar, avgrensa av ei kognitiv lyskjegle. Omsorga avgjer kva trykk designaren representerer. Forma som oppstår, viser kor langt omsorga rakk.

\subsection{Syntesegapet}

Kvar tradisjon gjev éin naudsynt komponent: formalgebraen gjev syntaks, C-K-teorien gjev epistemisk logikk, og multiskala-kompetansen gjev dynamikk. Ingen eksisterande rammeverk integrerer alle tre. Kjernalikninga i syntesen er:
\begin{equation*}
\mathrm{Form} = SG\!\left(\nabla_{C(A)}\, L(c,\, t)\right)
\end{equation*}
\noindent Form er grammatikken sin respons på landskapsgradienten slik agenten si lyskjegle oppfattar han.
